\documentclass[../../main.tex]{subfiles}% 注意这里的文件路径不能用 ./main.tex ,否则用latexmk编译子文件会报错
\graphicspath{{\subfix{./image/}}{\subfix{../../image/}}} % 指定图片引用路径,后续可以直接使用图片文件名
% 如果图片存放在上级目录,则使用 ../../image/ 作为备用路径

% 例如:
% \begin{figure}[H]
% \centering
% \includegraphics[width=0.3\textwidth]{图.png}
% \caption{}
% \label{figure:图}
% \end{figure}
% 注意:上述\label{}一定要放在\caption{}之后,否则引用图片序号会只会显示??.

\begin{document}
% SOURCE p=124; printed=109; object=ch03-notes-title; region=lower
\section{注记与参考文献}\label{section:ch03-notes-references}
% SOURCE p=124; printed=109; object=ch03-notes-lrn-history; region=bottom
\cref{section:ch03-lebesgue-radon-nikodym}：Lebesgue–Radon–Nikodym 定理由 Lebesgue\cite{folland-bib-092}在 $\mu$ 为 $\mathbb R^n$ 上 Lebesgue 测度的情形下证明。在假设 $\nu\ll\mu$ 下，Radon\cite{folland-bib-111}将其推广到 $\mathbb R^n$ 上任意正则 Borel 测度，Nikodym\cite{folland-bib-107}又将其推广到抽象空间上的测度。抽象情形下的 Lebesgue 分解见 Saks\cite{folland-bib-128}。正文中 Lebesgue–Radon–Nikodym 定理的证明与 Halmos\cite{folland-bib-061}中的证明相似，但更有效率；我从 L. Loomis 那里学到了这个证明。

% SOURCE p=125; printed=110; object=ch03-notes-variation-characterization; region=top
\cref{section:ch03-complex-measures}：
\begin{proposition}\label{proposition:ch03-notes-variation-partitions}
复测度 $\nu$ 的全变差有如下刻画：
\begin{align*}
|\nu|(E)=\sup\left\{\sum_{j=1}^n|\nu(E_j)|:
n\in\mathbb N,\ E_1,\ldots,E_n\text{ 两两不交},\ E=\bigcup_{j=1}^nE_j\right\}.
\end{align*}
\end{proposition}
通常把这一刻画（见\cref{exercise:ch03-complex-variation-suprema}）作为 $|\nu|$ 的定义。我们的定义似乎在更多场合有用，而且计算起来肯定更容易。
% SUPPLEMENT: id=v14-ch03-notes-variation-proof; parent=ch03-notes-variation-characterization; kind=proof+motivation+hypotheses; origin=independent-derivation
\begin{proof}
\textbf{有限分割刻画的补写证明（重构新增）。}这里 $E,E_j$ 都属于原测度空间的 $\sigma$-代数。任意有限不交分割满足
\begin{align*}
\sum_j|\nu(E_j)|\le\sum_j|\nu|(E_j)=|\nu|(E),
\end{align*}
故右端不超过左端。

反向，由极分解写 $\mathrm d\nu=h\,\mathrm d|\nu|$，其中 $|h|=1$ 几乎处处；在可测零测例外集上改取 $h=1$。给定 $0<\varepsilon<1$，把单位圆分为有限个不交 Borel 小弧 $A_1,\ldots,A_N$，并各选 $c_j\in A_j$，使 $|z-c_j|<\varepsilon$ 对 $z\in A_j$ 成立。令 $E_j=E\cap h^{-1}(A_j)$；它们是允许的有限分割，并且
\begin{align*}
|\nu(E_j)-c_j|\nu|(E_j)|
&\le\int_{E_j}|h-c_j|\,\mathrm d|\nu|
\le\varepsilon|\nu|(E_j),\\
\sum_j|\nu(E_j)|&\ge(1-\varepsilon)|\nu|(E).
\end{align*}
让 $\varepsilon\downarrow0$ 得反向不等式。空集可取单块空分割；若要求非空分块，删去空块即可。
\end{proof}
\begin{remark}
\textbf{动机与条件说明（重构新增）。}同一集合内的复质量可能相消；先分成小块再对各块质量取模，可逐渐消除这种抵消。可测性使各块的测度有定义，不交及穷尽 $E$ 保证质量既不重算也不遗漏。有限分割已足够，因为单位圆可分成有限个任意小的弧。有限复测度保证全变差有限，误差 $\varepsilon|\nu|(E)$ 因而趋零。此补证也给出了对应习题这一部分的独立证明，不以尚未写出的习题解答为前提。
\end{remark}
% END SUPPLEMENT: id=v14-ch03-notes-variation-proof

% SOURCE p=125; printed=110; object=ch03-notes-differentiation-history; region=upper-middle
\cref{section:ch03-euclidean-differentiation}：\cref{theorem:ch03-lebesgue-differentiation,theorem:ch03-measure-differentiation}由 Lebesgue\cite{folland-bib-092}给出，但我们采用的论证路线实质上是 Wiener\cite{folland-bib-161}的路线；一维情形的极大函数 $Hf$ 最早在 Hardy 与 Littlewood\cite{folland-bib-064}中得到研究。我们对\cref{theorem:ch03-ball-differentiation}的证明，展示了近年来被广泛运用的一种一般方法：通过适当极大函数的估计，控制一族算子的极限行为。

% SOURCE p=125; printed=110; object=ch03-notes-vitali-intro; region=middle
\cref{lemma:ch03-ball-covering}是 Wiener 覆盖引理的一个简化版本，取自 Rudin\cite{folland-bib-125}。另有一个年代更早、也更精细的覆盖定理，由 Vitali 给出，用于类似的目的：
% SOURCE p=125; printed=110; object=ch03-notes-vitali-covering; region=middle
\begin{theorem}[Vitali 覆盖定理]\label{theorem:ch03-vitali-covering}
若 $E\subseteq\mathbb R^n$，且 $\mathcal Q$ 是一个立方体族，使每个 $x\in E$ 都包含在 $\mathcal Q$ 中直径任意小的成员内，则存在一个有限或无限的两两不交序列 $\{Q_j\}\subseteq\mathcal Q$，使
\begin{align*}
m\left(E\setminus\bigcup_jQ_j\right)=0.
\end{align*}
\end{theorem}
证明可见许多著作，例如 Cohn\cite[\S6.2]{folland-bib-027}、Falconer\cite[\S1.3]{folland-bib-039}及 Hewitt 与 Stromberg\cite[\S17]{folland-bib-076}。
% SUPPLEMENT: id=v14-ch03-notes-vitali-proof; parent=ch03-notes-vitali-covering; kind=proof+motivation+hypotheses; origin=independent-derivation
\begin{remark}
\textbf{约定、动机与条件说明（重构新增）。}这里立方体具有正边长；若允许只含单点的退化立方体，取全部单点就会使结论失败。“任意小”意为对每个 $x\in E$ 和每个 $\delta>0$，都能选到包含 $x$ 且直径小于 $\delta$ 的成员。仅仅有一次普通覆盖不够：两个相交的单位闭区间覆盖其并，却无法选出不交子族覆盖到零测误差。立方体的形状把直径与体积统一联系起来，使下面的选取过程能用体积级数控制余集。证明先对闭立方体进行，再处理开闭边界约定；旋转立方体同样适用。没有预设 $E$ 可测：我们先证明余集外测度为零，再由 Lebesgue 完备性得到所写测度等式。
\end{remark}
\begin{proof}
\textbf{补写证明（重构新增）。}先设 $E\subset B=(0,1)^n$，只保留闭包包含在 $B$ 内的成员；每一点到 $B^c$ 的距离为正，所以这仍是任意细的覆盖。暂设成员为闭立方体。若候选族非空，选 $Q_1$，其直径大于候选直径上确界的一半；删除所有与 $Q_1$ 相交的候选。随后在剩余族中按同样方式选 $Q_2,Q_3,\ldots$，直到候选为空，或得到无限序列。每一步直径上确界为有限正数，故选取合法；所选立方体两两不交且包含在 $B$ 内。

若过程无限，则 $\sum_jm(Q_j)\le m(B)=1$。记 $d_j=\operatorname{diam}Q_j$；因 $m(Q_j)=(d_j/\sqrt n)^n$，有 $d_j\to0$。任一候选立方体 $Q$ 必在某个有限步骤被删除：否则它始终可选，便有 $d_j>\operatorname{diam}Q/2$ 对所有 $j$ 成立，矛盾。若 $j$ 是删除它的步骤，则 $Q\cap Q_j\ne\varnothing$，而在第 $j$ 步选取时 $Q$ 仍是候选，所以 $\operatorname{diam}Q<2d_j$。

设 $x\in E\setminus\bigcup_jQ_j$。固定 $N$，由于 $x$ 不属于有限个闭集 $Q_1,\ldots,Q_N$，可用任意细覆盖选一个包含 $x$ 且与这些闭集都不相交的立方体 $Q$。删除 $Q$ 的步骤必为 $j>N$。记 $c_j$ 为 $Q_j$ 的中心，利用一个交点，得
\begin{align*}
|x-c_j|\le\operatorname{diam}Q+\operatorname{diam}Q_j<3d_j.
\end{align*}
因此余集包含在 $\bigcup_{j>N}B(4d_j,c_j)$ 内。由球体积随半径的齐次性，存在仅依赖 $n$ 的常数 $C_n$，使 $m(B(4d_j,c_j))=C_nm(Q_j)$。外测度次可加性给出
\begin{align*}
m^*\left(E\setminus\bigcup_jQ_j\right)
\le C_n\sum_{j>N}m(Q_j)\longrightarrow0.
\end{align*}
若过程在有限步停止，未覆盖点会同样产生与全部已选闭集不交的新候选，矛盾，故此时已经覆盖 $E$。

对一般 $E\subseteq\mathbb R^n$，用开单位格子 $k+(0,1)^n$（$k\in\mathbb Z^n$）分别进行上述构造；每个格子内只选闭包仍在该格子中的候选。不同格子的选择互不相交。全部格子以外是可数个坐标超平面的并，其测度为零；这一点可先在有界盒子上用零厚度的长方体覆盖，再取可数并。把各格子的有限或可数序列合并，仍为可数不交序列；每格的余集以及格子边界的可数并均有外测度零，得到结论。

若原族采用开或半开立方体，先对其闭包族作上述选择，每个被选闭包选取一个对应原成员。闭包两两不交，故原成员也两两不交。原成员与其闭包的差只在边界上，立方体边界为有限个零测面之并；可数个被选边界的并仍为零测集。因此改回原成员不影响结论。
\end{proof}
% END SUPPLEMENT: id=v14-ch03-notes-vitali-proof

% SOURCE p=125; printed=110; object=ch03-notes-bv-history; region=middle-lower
\cref{section:ch03-bounded-variation}：本节的主要结果由 Lebesgue 和 Vitali 给出，详细文献可见 Hawkins\cite{folland-bib-070}。\cref{exercise:ch03-stieltjes-change-variables}给出了 Lebesgue 积分换元公式的一种形式；其他形式可见 Serrin 与 Varberg\cite{folland-bib-133}。\cref{exercise:ch03-monotone-series-derivative}是 Fubini 的一个定理，而\cref{exercise:ch03-strict-singular-function}中的例子由 Brown\cite{folland-bib-021}给出。

% SOURCE p=125; printed=110; object=ch03-notes-riemann-stieltjes; region=lower
\begin{definition}[Riemann–Stieltjes 积分的原始构想]\label{definition:ch03-riemann-stieltjes-integral}
Stieltjes 积分 $\int_a^b g\,\mathrm dF$ 最初是在 $F$ 为 $[a,b]$ 上递增函数的假设下，定义为 Riemann 和 $\sum g(t_j)[F(t_j)-F(t_{j-1})]$ 的极限。
\end{definition}
这种“Riemann–Stieltjes”积分的理论与普通 Riemann 积分非常相似，但在允许 $g,F$ 都不连续时，需要小心处理。参见 ter Horst\cite{folland-bib-148}，其中包含\cref{theorem:ch02-riemann-lebesgue}对 Stieltjes 积分的对应结果。
% SUPPLEMENT: id=v14-ch03-notes-riemann-stieltjes-remark; parent=ch03-notes-riemann-stieltjes; kind=motivation+hypotheses+scope; origin=independent-derivation
\begin{remark}
\textbf{动机与条件说明（重构新增）。}每段增量 $F(t_j)-F(t_{j-1})$ 充当长度权重；$F$ 递增时权重非负，且总和为 $F(b)-F(a)$。原文在此只回顾构想，并未规定取样点、趋细方式或端点约定，因此这一段不应被误用为涵盖所有不连续函数的唯一完整定义。若 $g$ 连续且 $F$ 连续递增，可取分割 $a=x_0<\cdots<x_N=b$、任意标记 $\xi_j\in[x_{j-1},x_j]$ 和网格宽度趋零；所得极限与 Lebesgue–Stieltjes 积分相等，理由如下。
\end{remark}
\begin{proof}
\textbf{连续情形的补写验证（重构新增）。}用常值延拓后的 $F$ 定义区间 Stieltjes 测度 $\mu_F$；连续性保证端点无原子。在 $(x_{j-1},x_j]$ 上令 $s_P=g(\xi_j)$。则 $\int s_P\,\mathrm d\mu_F$ 正是带标记的 Stieltjes 和。紧区间上 $g$ 一致连续，故其连续模 $\omega_g(\delta)$ 随 $\delta\downarrow0$ 趋零，而
\begin{align*}
\left|\int s_P\,\mathrm d\mu_F-\int g\,\mathrm d\mu_F\right|
\le\omega_g(|P|)\,[F(b)-F(a)]\longrightarrow0.
\end{align*}
这也说明标记选择不影响极限。若 $F$ 恒定，总质量为零，两者同为零。不连续情形需另定约定，不能从这一验证直接推广。
\end{proof}
% END SUPPLEMENT: id=v14-ch03-notes-riemann-stieltjes-remark

% SOURCE p=125; printed=110; object=ch03-notes-countable-exceptions-intro; region=lower
Cantor 函数表明，一个连续且几乎处处可微的函数不必是其导数的积分。下面是一个很不平凡的定理。
% SOURCE p=125; printed=110; object=ch03-notes-countable-exceptions; region=lower
\begin{theorem}\label{theorem:ch03-countable-exception-ftc}
若 $F$ 在 $[a,b]$ 上连续，$F'(x)$ 对每个 $x\in[a,b]\setminus A$ 都存在，其中 $A$ 可数，并且 $F'\in L^1$，则 $F$ 绝对连续，从而可由对 $F'$ 积分而恢复。
\end{theorem}
证明可见 Cohn\cite[\S6.3]{folland-bib-027}；对于较容易的 $A=\varnothing$ 情形，也可参见 Rudin\cite[定理 7.26]{folland-bib-125}。
% SUPPLEMENT: id=v14-ch03-notes-countable-exceptions-proof; parent=ch03-notes-countable-exceptions; kind=proof+motivation+hypotheses; origin=independent-derivation
\begin{remark}
\textbf{动机与条件说明（重构新增）。}这里 $a<b$ 为有限端点，$F$ 可为实值或复值；端点可放进可数例外集，因而不需另作端点双侧导数约定。$F'$ 在例外集上任取有限值不改变 $L^1$ 类。可数性比零测性更强，允许给每个异常点分配一个可求和误差；若只假定异常集零测，Cantor 函数已是否定例。连续性用于异常点处控制函数增量，不能删去：有跳跃的阶跃函数在一个点外可微，导数几乎处处为零，却不绝对连续。下面利用\cref{section:ch02-notes-references}已经定义的尺度积分，并把所需的桥接结论一并证明。
\end{remark}
\begin{proof}
\textbf{补写证明（重构新增）。}我们先补齐三个步骤，再应用到 $F$。记带标记分割的每段为 $I_j=[x_{j-1},x_j]$，标记为 $t_j\in I_j$，长度为 $\ell_j$；尺度条件为 $\ell_j<\delta(t_j)$。

\textbf{第一步：细分割存在且积分值唯一。}给定任意处处正的尺度 $\delta$。若 $[a,b]$ 不存在 $\delta$-细分割，将它二等分，至少一半也没有这样的分割，否则把两半的分割拼起来即可。不断二分，得到嵌套闭区间 $I_k$，长度趋零且每个都没有细分割。紧致性给出唯一公共点 $t$。当 $|I_k|<\delta(t)$ 时，以 $t$ 为唯一标记的单段分割就是细分割，矛盾。积分值唯一性则取两个候选值相应尺度的逐点最小值；存在的共同细分割使两候选值之差小于任意两倍误差，故相等。

\textbf{第二步：有限值的 Lebesgue 可积函数有相同的尺度积分。}先设 $f:[a,b]\to\mathbb R$ 处处有限且可积。固定 $\eta>0$，按值域分层
\begin{align*}
E_k=\{x:k\eta\le f(x)<(k+1)\eta\},\qquad k\in\mathbb Z.
\end{align*}
令 $\rho(H)=\int_H(1+|f|)\,\mathrm dm$。给定 $\gamma>0$，对每个 $k$ 选相对开集 $U_k\supseteq E_k$，使
\begin{align*}
\rho(U_k\setminus E_k)<\frac{\gamma w_k}{1+|k|\eta},
\qquad w_k=\frac{2^{-|k|}}3,\qquad\sum_{k\in\mathbb Z}w_k=1.
\end{align*}
这种选择来自 Lebesgue 外正则性与 $|f|$ 积分的绝对连续性：先使开集多出的长度足够小，便同时控制长度和 $|f|$ 积分。对 $t\in E_k$ 取 $\delta(t)>0$，使 $[a,b]\cap(t-\delta(t),t+\delta(t))\subseteq U_k$。

给定任一细分割，把标记属于 $E_k$ 的各段的并记为 $B_k$；忽略有限个端点，所有段及各 $B_k$ 两两不交，且 $B_k\subseteq U_k$。在 $I_j$ 的标记为 $t_j\in E_k$ 时，$|f(t_j)-k\eta|<\eta$。在 $B_k\cap E_k$ 上也有 $|f-k\eta|\le\eta$，而在余下部分用 $|f-k\eta|\le|f|+|k|\eta$。于是
\begin{align*}
\left|\sum_jf(t_j)\ell_j-\int_a^bf\,\mathrm dm\right|
&\le\sum_j\int_{I_j}|f(t_j)-f(x)|\,\mathrm dx\\
&\le2\eta(b-a)+\sum_k(1+|k|\eta)\rho(U_k\setminus E_k)\\
&<2\eta(b-a)+\gamma.
\end{align*}
依次取足够小的 $\eta,\gamma$ 即得尺度积分等于 Lebesgue 积分。复值情形对实部和虚部分别应用此结论，再取两尺度的最小值；其线性组合由和的三角不等式满足同一定义。若原 $L^1$ 代表只几乎处处有限，在零测例外上先赋有限值，得到本步骤需要的代表。

\textbf{第三步：连续函数在可数集外的导数给出尺度积分。}不重复枚举 $A\cup\{a,b\}=\{c_k\}$，有限情形使用有限枚举。在该集合外令 $f=F'$，在集合上令 $f=0$。固定 $\varepsilon>0$。若 $t\notin A\cup\{a,b\}$，可微性允许选择 $\delta(t)>0$，使
\begin{align*}
|F(y)-F(t)-f(t)(y-t)|\le\varepsilon|y-t|
\quad\text{当 } |y-t|<\delta(t),\ y\in[a,b].
\end{align*}
若 $t=c_k$，由连续性选 $\delta(c_k)>0$，使相同邻域内
\begin{align*}
|F(y)-F(c_k)|<\varepsilon 2^{-k-2}.
\end{align*}
对细分割中的正常标记，把端点 $x_{j-1},x_j$ 分别代入可微性估计，并相减，误差至多为 $\varepsilon\ell_j$。对异常标记 $c_k$，两端函数值差的模小于 $\varepsilon2^{-k-1}$；一个固定点在分割中至多能作为两段的标记，因此所有异常标记的误差总和小于 $\varepsilon\sum_k2^{-k}\le\varepsilon$。将各段相加、望远镜相消，得
\begin{align*}
\left|F(b)-F(a)-\sum_jf(t_j)\ell_j\right|
\le\varepsilon(b-a)+\varepsilon.
\end{align*}
这正证明 $f$ 的尺度积分存在且等于 $F(b)-F(a)$。此步骤只用可微性和异常点处连续性，并不假定 $f\in L^1$。

现在回到定理，题设使第三步中的 $f$ 属于 $L^1$。第二步和积分值唯一性给出 $F(b)-F(a)=\int_a^bf\,\mathrm dm$。对每个子区间 $[a,x]$ 重复同一论证，得到 $F(x)-F(a)=\int_a^xf\,\mathrm dm$；$x=a$ 的等式显然。由已证明的 Lebesgue 微积分基本定理，$F$ 绝对连续，并由其导数的积分逐点恢复。
\end{proof}
% END SUPPLEMENT: id=v14-ch03-notes-countable-exceptions-proof

% SOURCE p=125; printed=110; object=ch03-notes-nonintegrable-derivative; region=bottom
然而，事情并未到此结束。
\begin{proposition}\label{proposition:ch03-nonintegrable-derivative-existence}
存在处处可微的函数 $F$，使 $F'\notin L^1$。
\end{proposition}
% SOURCE p=125; printed=110; object=ch03-notes-nonintegrable-example; region=bottom
\begin{example}\label{example:ch03-nonintegrable-derivative}
或许最简单的例子是 $F(x)=x^2\sin(x^{-2})$（见\cref{exercise:ch03-differentiable-non-bv}）。
\end{example}
% SOURCE p=125,126; printed=110,111; object=ch03-notes-improper-derivative-integral; region=bottom-top
这里的问题只在 $x=0$，所以对于 $a\le0\le b$，可以把 $\int_a^bF'(t)\,\mathrm dt$ 看作反常积分，即当 $\varepsilon\to0$ 时，在 $[a,b]\setminus[-\varepsilon,\varepsilon]$ 上的 Lebesgue 积分之极限。
% SUPPLEMENT: id=v14-ch03-notes-nonintegrable-example-proof; parent=ch03-notes-improper-derivative-integral; kind=proof+motivation+hypotheses; origin=independent-derivation
\begin{proof}
\textbf{不可绝对积分及反常恢复的补写验证（重构新增）。}按所引习题，在零点令 $F(0)=0$，则 $F'(0)=0$，而对 $x\ne0$，
\begin{align*}
F'(x)=2x\sin(x^{-2})-2x^{-1}\cos(x^{-2}).
\end{align*}
第一项在有界区间上绝对可积。对任意 $\eta>0$，代换 $u=x^{-2}$ 给出
\begin{align*}
\int_0^\eta \frac{2|\cos(x^{-2})|}{x}\,\mathrm dx
=\int_{\eta^{-2}}^\infty\frac{|\cos u|}{u}\,\mathrm du=\infty.
\end{align*}
最后的发散可由足够大的整数 $k$ 对应区间 $[2\pi k-\pi/3,2\pi k+\pi/3]$ 看出：其上 $|\cos u|\ge1/2$，积分至少为一个固定正常数乘 $1/k$，而这些区间不交。反三角不等式说明 $\int_0^\eta|F'|=\infty$。负半轴由 $F$ 为偶函数而 $F'$ 为奇函数得到同样结论。

另一方面，$F$ 在避开零点的紧区间上光滑。因此当 $a<0<b$ 且 $\varepsilon$ 足够小时，经典基本定理给出
\begin{align*}
\int_a^{-\varepsilon}F'(t)\,\mathrm dt+\int_\varepsilon^bF'(t)\,\mathrm dt
=F(-\varepsilon)-F(a)+F(b)-F(\varepsilon)
\longrightarrow F(b)-F(a).
\end{align*}
当 $a=0$ 或 $b=0$ 时省去相应空侧，仍由 $F(\pm\varepsilon)\to F(0)$ 得同一极限；$a=b=0$ 时约定积分为零。左右两侧反常积分各自收敛，所以这里不只是对称主值的抵消。
\end{proof}
\begin{remark}
\textbf{动机与条件说明（重构新增）。}振荡使导数的正负贡献能够条件抵消，而其绝对值积分发散。这足以说明处处可微不保证给定区间上的导数属于 $L^1$。把零点取值补为零是连续性与可微性的必要约定；离开零点的光滑性支持经典积分计算，零点连续性则使反常端点项有极限。
\end{remark}
% END SUPPLEMENT: id=v14-ch03-notes-nonintegrable-example-proof

% SOURCE p=126; printed=111; object=ch03-notes-nowhere-integrable-claim; region=top
不过，不难构造这样的例子：$F'$ 的奇点如此复杂，以至于 $F'$ 在任何区间上都不 Lebesgue 可积。在这种情形下，Lebesgue 积分确实不够用。
% SUPPLEMENT: id=v14-ch03-notes-nowhere-claim-qualification; parent=ch03-notes-nowhere-integrable-claim; kind=source-qualification+proof+motivation+hypotheses; origin=independent-derivation
\begin{remark}
\textbf{前提范围澄清（重构新增）。}若此句继续沿用上一段“$F$ 处处可微且导数有限”的前提，那么“在任何区间上都不可积”不能成立。下面给出直接证明：在任意非退化开区间内，都存在一个更小的非退化闭区间，使 $F'$ 有界且 Lebesgue 可积。因此不能为这一字面组合补造例子。与上下文及后面基本定理相容的较弱表述是：处处可微的函数，其导数在给定的整个闭区间上仍可能不属于 $L^1$；前一例已经完整证明这一点。原文句子在上方保留，此处明确区分原述与重构澄清。
\end{remark}
\begin{proof}
\textbf{局部有界子区间的证明（重构新增）。}设 $F$ 在某开区间 $I$ 上处处可微，任取非退化闭区间 $J=[c,d]\subset I$。选择 $n_0$，使 $x+1/n\in I$ 对所有 $x\in J,n\ge n_0$ 成立。函数
\begin{align*}
D_n(x)=n\,[F(x+1/n)-F(x)]
\end{align*}
在 $J$ 上连续，并逐点趋于有限值 $F'(x)$。故对每个 $x$，序列 $\{D_n(x):n\ge n_0\}$ 有界。令
\begin{align*}
C_k=\bigcap_{n\ge n_0}\{x\in J:|D_n(x)|\le k\},\qquad k\in\mathbb N.
\end{align*}
这些集合在 $J$ 中闭，且 $J=\bigcup_kC_k$。至少一个 $C_k$ 在 $J$ 中有非空相对内部：否则可递归选取非退化闭区间 $J_k\subseteq\operatorname{int}J_{k-1}\setminus C_k$，使长度小于 $2^{-k}$，从 $J_0=J$ 开始。闭性与空内部保证每一步有一个开子区间可选。嵌套紧区间的公共点将不在任何 $C_k$ 中，与它们覆盖 $J$ 矛盾。

因此存在非退化闭子区间 $K\subseteq C_k$。令 $n\to\infty$，得 $|F'|\le k$ 在 $K$ 上成立。又 $F'$ 是连续差商的逐点极限，故 Borel 可测；$m(K)<\infty$，所以 $F'\in L^1(K)$。原 $I$ 可以取在任意给定非退化区间内部，这就排除了上述处处可微前提下的“任何区间都不可积”。复值情形用复模的闭性与实部虚部可测性同样成立。
\end{proof}
% END SUPPLEMENT: id=v14-ch03-notes-nowhere-claim-qualification

% SOURCE p=126; printed=111; object=ch03-notes-general-ftc-intro; region=top
然而，\cref{section:ch02-notes-references}讨论过的 Henstock–Kurzweil 积分（或 Denjoy 积分、Perron 积分）足以对这样的 $F'$ 积分；使用这种积分便得到一般的微积分基本定理：
% SOURCE p=126; printed=111; object=ch03-notes-general-ftc; region=top
\begin{theorem}[一般的微积分基本定理]\label{theorem:ch03-hk-fundamental-calculus}
若 $F$ 在 $[a,b]$ 上处处可微，则
\begin{align*}
F(b)-F(a)=\int_a^bF'(t)\,\mathrm dt,
\end{align*}
其中积分按 Henstock–Kurzweil 意义理解。
\end{theorem}
% SUPPLEMENT: id=v14-ch03-notes-general-ftc-proof; parent=ch03-notes-general-ftc; kind=proof+motivation+hypotheses; origin=independent-derivation
\begin{proof}
\textbf{补写证明（重构新增）。}使用\cref{theorem:ch03-countable-exception-ftc}补证中的第三步，取内部例外集为空，即由每个标记处的线性近似控制逐段误差，得到导数的尺度积分等于端点差；那一步没有使用 $F'\in L^1$。若在 $a,b$ 给定单侧导数，也可直接在那里使用单侧线性近似。或者仍将两个端点归入异常点，以零值代替端点导数；这不改变尺度积分：只在有限个点取非零有限值的函数，对每个该点把允许区间长度取得足够小，且每点至多被用作两段标记，就使所有带标记和的绝对值任意小。尺度积分的线性性由取尺度最小值并相加得到，故改回真实端点导数后积分值相同。
\end{proof}
\begin{remark}
\textbf{动机与条件说明（重构新增）。}尺度可依赖标记点，使局部线性近似在导数变化很剧烈时仍逐段有效；逐段误差按长度求和，最终只留下端点差。处处可微自动给出连续性，却不要求导数连续、有界或绝对可积。这里的积分类型必须明确，不能把结论误读为任意导数均 Lebesgue 可积；$x^2\sin(x^{-2})$ 的例子正好说明区别。
\end{remark}
% END SUPPLEMENT: id=v14-ch03-notes-general-ftc-proof
\end{document}
